\begin{proposition}
\label{proposition-etale-locally-standard}
Let $R \to S$ be a ring map. Let $\mathfrak q \subset S$ be a prime.
If $R \to S$ is \'etale at $\mathfrak q$, then there exists
a $g \in S$, $g \not \in \mathfrak q$ such that $R \to S_g$
is standard \'etale.
\end{proposition}

\begin{proof}
The following proof is a little roundabout and there may be ways to
shorten it.

\medskip\noindent
Step 1. By Definition \ref{definition-etale}
there exists a $g \in S$, $g \not \in \mathfrak q$
such that $R \to S_g$ is \'etale. Thus we may assume that $S$ is \'etale
over $R$.

\medskip\noindent
Step 2. By Lemma \ref{lemma-etale} there exists an \'etale ring map
$R_0 \to S_0$ with $R_0$ of finite type over $\mathbf{Z}$, and a ring map
$R_0 \to R$ such that $R = R \otimes_{R_0} S_0$. Denote
$\mathfrak q_0$ the prime of $S_0$ corresponding to $\mathfrak q$.
If we show the result for $(R_0 \to S_0, \mathfrak q_0)$ then the
result follows for $(R \to S, \mathfrak q)$ by base change. Hence
we may assume that $R$ is Noetherian.

\medskip\noindent
Step 3.
Note that $R \to S$ is quasi-finite by Lemma \ref{lemma-etale-quasi-finite}.
By Lemma \ref{lemma-quasi-finite-open-integral-closure}
there exists a finite ring map $R \to S'$, an $R$-algebra map
$S' \to S$, an element $g' \in S'$ such that
$g' \not \in \mathfrak q$ such that $S' \to S$ induces
an isomorphism $S'_{g'} \cong S_{g'}$.
(Note that of course $S'$ is not \'etale over $R$ in general.)
Thus we may assume that (a) $R$ is Noetherian, (b) $R \to S$ is finite
and (c) $R \to S$ is \'etale at $\mathfrak q$
(but no longer necessarily \'etale at all primes).

\medskip\noindent
Step 4. Let $\mathfrak p \subset R$ be the prime corresponding
to $\mathfrak q$. Consider the fibre ring
$S \otimes_R \kappa(\mathfrak p)$. This is a finite algebra over
$\kappa(\mathfrak p)$. Hence it is Artinian
(see Lemma \ref{lemma-finite-dimensional-algebra}) and
so a finite product of local rings
$$
S \otimes_R \kappa(\mathfrak p) = \prod\nolimits_{i = 1}^n A_i
$$
see Proposition \ref{proposition-dimension-zero-ring}. One of the factors,
say $A_1$, is the local ring $S_{\mathfrak q}/\mathfrak pS_{\mathfrak q}$
which is isomorphic to $\kappa(\mathfrak q)$,
see Lemma \ref{lemma-etale-at-prime}. The other factors correspond to
the other primes, say $\mathfrak q_2, \ldots, \mathfrak q_n$ of
$S$ lying over $\mathfrak p$.

\medskip\noindent
Step 5. We may choose a nonzero element $\alpha \in \kappa(\mathfrak q)$ which
generates the finite separable field extension
$\kappa(\mathfrak q)/\kappa(\mathfrak p)$ (so even if the
field extension is trivial we do not allow $\alpha = 0$).
Note that for any $\lambda \in \kappa(\mathfrak p)^*$ the
element $\lambda \alpha$ also generates $\kappa(\mathfrak q)$
over $\kappa(\mathfrak p)$. Consider the element
$$
\overline{t} =
(\alpha, 0, \ldots, 0) \in
\prod\nolimits_{i = 1}^n A_i =
S \otimes_R \kappa(\mathfrak p).
$$
After possibly replacing $\alpha$ by $\lambda \alpha$ as above
we may assume that $\overline{t}$ is the image of $t \in S$.
Let $I \subset R[x]$ be the kernel of the $R$-algebra
map $R[x] \to S$ which maps $x$ to $t$. Set $S' = R[x]/I$,
so $S' \subset S$. Here is a diagram
$$
\xymatrix{
R[x] \ar[r] & S' \ar[r] & S \\
R \ar[u] \ar[ru] \ar[rru] & &
}
$$
By construction the primes $\mathfrak q_j$, $j \geq 2$ of $S$ all
lie over the prime $(\mathfrak p, x)$ of $R[x]$, whereas
the prime $\mathfrak q$ lies over a different prime of $R[x]$
because $\alpha \not = 0$.

\medskip\noindent
Step 6. Denote $\mathfrak q' \subset S'$ the prime of $S'$
corresponding to $\mathfrak q$. By the above $\mathfrak q$ is
the only prime of $S$ lying over $\mathfrak q'$. Thus we see that
$S_{\mathfrak q} = S_{\mathfrak q'}$, see
Lemma \ref{lemma-unique-prime-over-localize-below} (we have
going up for $S' \to S$ by Lemma \ref{lemma-integral-going-up}
since $S' \to S$ is finite as $R \to S$ is finite).
It follows that $S'_{\mathfrak q'} \to S_{\mathfrak q}$ is finite
and injective as the localization of the finite injective ring map
$S' \to S$. Consider the maps of local rings
$$
R_{\mathfrak p} \to S'_{\mathfrak q'} \to S_{\mathfrak q}
$$
The second map is finite and injective. We have
$S_{\mathfrak q}/\mathfrak pS_{\mathfrak q} = \kappa(\mathfrak q)$,
see Lemma \ref{lemma-etale-at-prime}.
Hence a fortiori
$S_{\mathfrak q}/\mathfrak q'S_{\mathfrak q} = \kappa(\mathfrak q)$.
Since
$$
\kappa(\mathfrak p) \subset \kappa(\mathfrak q') \subset \kappa(\mathfrak q)
$$
and since $\alpha$ is in the image of $\kappa(\mathfrak q')$ in
$\kappa(\mathfrak q)$
we conclude that $\kappa(\mathfrak q') = \kappa(\mathfrak q)$.
Hence by Nakayama's Lemma \ref{lemma-NAK} applied to the
$S'_{\mathfrak q'}$-module map $S'_{\mathfrak q'} \to S_{\mathfrak q}$,
the map $S'_{\mathfrak q'} \to S_{\mathfrak q}$ is surjective.
In other words,
$S'_{\mathfrak q'} \cong S_{\mathfrak q}$.

\medskip\noindent
Step 7. By Lemma \ref{lemma-isomorphic-local-rings} there exist
$g \in S$, $g \not \in \mathfrak q$ and $g' \in S'$, $g' \not \in \mathfrak q'$
such that $S'_{g'} \cong S_g$. As $R$ is Noetherian the ring $S'$ is finite
over $R$ because it is an $R$-submodule
of the finite $R$-module $S$. Hence after replacing $S$ by $S'$ we may
assume that (a) $R$ is Noetherian, (b) $S$ finite over $R$, (c)
$S$ is \'etale over $R$ at $\mathfrak q$, and (d) $S = R[x]/I$.

\medskip\noindent
Step 8. Consider the ring
$S \otimes_R \kappa(\mathfrak p) = \kappa(\mathfrak p)[x]/\overline{I}$
where $\overline{I} = I \cdot \kappa(\mathfrak p)[x]$ is the ideal generated
by $I$ in $\kappa(\mathfrak p)[x]$. As $\kappa(\mathfrak p)[x]$ is a PID
we know that $\overline{I} = (\overline{h})$ for some monic
$\overline{h} \in \kappa(\mathfrak p)[x]$. After replacing $\overline{h}$
by $\lambda \cdot \overline{h}$ for some $\lambda \in \kappa(\mathfrak p)$
we may assume that $\overline{h}$ is the image of some $h \in I \subset R[x]$.
(The problem is that we do not know if we may choose $h$ monic.)
Also, as in Step 4 we know that
$S \otimes_R \kappa(\mathfrak p) = A_1 \times \ldots \times A_n$ with
$A_1 = \kappa(\mathfrak q)$ a finite separable extension of
$\kappa(\mathfrak p)$ and $A_2, \ldots, A_n$ local. This implies
that
$$
\overline{h} = \overline{h}_1 \overline{h}_2^{e_2} \ldots \overline{h}_n^{e_n}
$$
for certain pairwise coprime irreducible monic polynomials
$\overline{h}_i \in \kappa(\mathfrak p)[x]$ and certain
$e_2, \ldots, e_n \geq 1$. Here the numbering is chosen so that
$A_i = \kappa(\mathfrak p)[x]/(\overline{h}_i^{e_i})$ as
$\kappa(\mathfrak p)[x]$-algebras. Note that $\overline{h}_1$ is
the minimal polynomial of $\alpha \in \kappa(\mathfrak q)$ and hence
is a separable polynomial (its derivative is prime to itself).

\medskip\noindent
Step 9. Let $m \in I$ be a monic element; such an element exists
because the ring extension $R \to R[x]/I$ is finite hence integral.
Denote $\overline{m}$ the image in $\kappa(\mathfrak p)[x]$.
We may factor
$$
\overline{m} = \overline{k}
\overline{h}_1^{d_1} \overline{h}_2^{d_2} \ldots \overline{h}_n^{d_n}
$$
for some $d_1 \geq 1$, $d_j \geq e_j$, $j = 2, \ldots, n$ and
$\overline{k} \in \kappa(\mathfrak p)[x]$ prime to all the $\overline{h}_i$.
Set $f = m^l + h$ where $l \deg(m) > \deg(h)$, and $l \geq 2$.
Then $f$ is monic as a polynomial over $R$. Also, the image $\overline{f}$
of $f$ in $\kappa(\mathfrak p)[x]$ factors as
$$
\overline{f} =
\overline{h}_1 \overline{h}_2^{e_2} \ldots \overline{h}_n^{e_n}
+
\overline{k}^l \overline{h}_1^{ld_1} \overline{h}_2^{ld_2}
\ldots \overline{h}_n^{ld_n}
=
\overline{h}_1(\overline{h}_2^{e_2} \ldots \overline{h}_n^{e_n}
+
\overline{k}^l
\overline{h}_1^{ld_1 - 1} \overline{h}_2^{ld_2} \ldots \overline{h}_n^{ld_n})
= \overline{h}_1 \overline{w}
$$
with $\overline{w}$ a polynomial relatively prime to $\overline{h}_1$.
Set $g = f'$ (the derivative with respect to $x$).

\medskip\noindent
Step 10. The ring map $R[x] \to S = R[x]/I$ has the properties:
(1) it maps $f$ to zero, and
(2) it maps $g$ to an element of $S \setminus \mathfrak q$.
The first assertion is clear since $f$ is an element of $I$.
For the second assertion we just have to show that $g$ does
not map to zero in
$\kappa(\mathfrak q) = \kappa(\mathfrak p)[x]/(\overline{h}_1)$.
The image of $g$ in $\kappa(\mathfrak p)[x]$ is the derivative
of $\overline{f}$. Thus (2) is clear because
$$
\overline{g} =
\frac{\text{d}\overline{f}}{\text{d}x} =
\overline{w}\frac{\text{d}\overline{h}_1}{\text{d}x} +
\overline{h}_1\frac{\text{d}\overline{w}}{\text{d}x},
$$
$\overline{w}$ is prime to $\overline{h}_1$ and
$\overline{h}_1$ is separable.

\medskip\noindent
Step 11.
We conclude that $\varphi : R[x]/(f) \to S$ is a surjective ring map,
$R[x]_g/(f)$ is \'etale over $R$ (because it is standard \'etale,
see Lemma \ref{lemma-standard-etale}) and $\varphi(g) \not \in \mathfrak q$.
Pick an element $g' \in R[x]/(f)$ such that
also $\varphi(g') \not \in \mathfrak q$ and $S_{\varphi(g')}$
is \'etale over $R$ (which exists since $S$ is \'etale over $R$ at
$\mathfrak q$). Then the ring map
$R[x]_{gg'}/(f) \to S_{\varphi(gg')}$ is a surjective map of \'etale
algebras over $R$. Hence it is \'etale by Lemma \ref{lemma-map-between-etale}.
Hence it is a localization by
Lemma \ref{lemma-surjective-flat-finitely-presented}.
Thus a localization of $S$ at an element not in $\mathfrak q$ is
isomorphic to a localization of a standard \'etale algebra over $R$
which is what we wanted to show.
\end{proof}